% ------------------------------------------------------------
% §5  Iterating the function: the dynamics of SHA composed with itself
% ------------------------------------------------------------
\section{Iterating the function}\label{sec:iterated}

\Cref{sec:randommaps} pointed Flajolet's telescope at a single
application of \SHA{}. But the function is very often applied to its own
output---%
\[
  \varphi(x) \;=\; \SHA(x), \qquad
  \varphi^{k} \;=\; \underbrace{\SHA \circ \SHA \circ \cdots \circ
  \SHA}_{k},
\]
read as a self-map of the $256$-bit strings (feed the digest back as a
one-block message). This is no idle variation. Password stretching runs
$\varphi^{k}$ for $k$ in the hundreds of thousands; hash chains for
one-time passwords and hash-based signatures are iteration made into a
protocol; Bitcoin's core primitive is literally $\varphi^{2}$, the
double hash $\SHA(\SHA(\cdot))$. And iteration is its own coordinate
system for the random-function conjecture: the \emph{dynamics} of
$\varphi$---how its image shrinks, where its orbits settle, whether it
hides fixed points---are a fresh battery of falsifiable predictions,
sharper in some ways than the static silhouette of \cref{sec:randommaps}
because they probe the function's interaction with \emph{itself}. This
section reads those predictions off the random-map model and asks, in
each case, what an anomaly would mean. The user's instinct is the right
place to start: one might imagine iteration eventually shuffling the
state space like a random permutation---and seeing \emph{exactly why
that almost happens, and exactly what stops it}, turns out to expose the
single design decision that makes \SHA{} one-way.

\subsection{Iteration leaks: the shrinking image}\label{sec:shrinking}

A permutation loses nothing: applied repeatedly, it merely relabels. A
random-looking map is different, and the difference compounds. Recall
from \cref{sec:fungraph} that one application of a random map covers only
$(1-e^{-1})N \approx 0.63\,N$ of the space; the leaves---the $37\%$ of
points with no preimage---are gone after one step and never return.
Apply $\varphi$ again and the survivors thin again, by the same logic
applied to the image. Writing $\beta_{k}$ for the expected fraction of
the space still reachable after $k$ steps, the recurrence is exactly
\[
  \beta_{0} = 1, \qquad \beta_{k+1} = 1 - e^{-\beta_{k}},
\]
(``a point survives to step $k+1$ iff at least one of its
$\varphi$-preimages survived to step $k$''), and its solution decays
with a clean asymptotic law~\citep{FlajoletOdlyzko1990}:
\[
  \beta_{k} \;\sim\; \frac{2}{k} \qquad (k \to \infty).
\]
Iteration is a funnel (\cref{fig:funnel}). After $k$ rounds the reachable
image has shrunk to about $2N/k$ points, so the process has quietly
discarded roughly $\log_{2}(k/2)$ bits of range. The number is small but
not zero and the moral is exact: \emph{composing a hash with itself can
only lose entropy, never create it}. For a password-stretching count of
$k = 10^{6}$, the naive iterate $\varphi^{k}$ confines its outputs to a
$2^{-19}$ fraction of the digest space---nineteen bits of the advertised
$256$ silently surrendered to the funnel.

\begin{figure}[ht]
  \centering
  \begin{tikzpicture}[font=\small, >={Stealth[length=1.6mm]}]
    % funnel: nested horizontal bars shrinking downward
    \foreach \y/\w/\lab in {0/4.4/{$\State$: all $N=2^{256}$},
        -0.72/2.78/{$\varphi(\State)\approx 0.63\,N$},
        -1.34/2.06/{$\varphi^{2}(\State)\approx 0.47\,N$},
        -1.86/1.30/{$\varphi^{4}(\State)\approx 0.37\,N$}} {
      \draw[fill=projfill, draw=projframe]
        (-\w/2,\y) rectangle (\w/2,\y-0.42);
    }
    \node[font=\scriptsize] at (0,-0.21) {$N$};
    \node[right, font=\scriptsize] at (2.3,-0.21) {all strings};
    \node[right, font=\scriptsize] at (2.3,-0.93) {leaves gone (37\%)};
    \node[right, font=\scriptsize] at (2.3,-2.07)
      {$\approx 2N/k$ after $k$ steps};
    % dots
    \node at (0,-2.42) {$\vdots$};
    % recurrent set at the bottom
    \draw[fill=lenfill, draw=orange!70!black]
      (-0.62,-3.2) rectangle (0.62,-2.85);
    \node[font=\scriptsize] at (0,-3.02) {$R$};
    \node[right, font=\scriptsize, align=left] at (2.3,-3.02)
      {recurrent set $R=\bigcap_{k}\varphi^{k}(\State)$:\\
       the cycles, $\approx\!\sqrt{\pi N/2}=2^{128.3}$ points};
    \draw[->] (0,-2.5) -- (0,-2.83);
    % downward funnel guide lines
    \draw[gray, dashed] (-2.2,0) -- (-0.62,-2.85);
    \draw[gray, dashed] (2.2,0) -- (0.62,-2.85);
    \node[left, font=\scriptsize, align=right] at (-2.3,-1.5)
      {\emph{entropy}\\ \emph{funnel}\\ $\beta_{k}\!\sim\!2/k$};
  \end{tikzpicture}
  \caption{Iterating a random-looking map is a one-way funnel. Each
  application discards the current leaves, so the reachable image shrinks
  as $\beta_{k}\sim 2/k$ toward the \emph{recurrent set} $R$---the union
  of all cycles, about $\sqrt{\pi N/2}\approx 2^{128.3}$ points. On $R$,
  and only there, $\varphi$ acts as a permutation (\cref{sec:permutation}).
  The funnel is why iterated hashing loses (a little) entropy, and why
  real key-stretching functions are \emph{not} pure iteration.}
  \label{fig:funnel}
\end{figure}

Two consequences deserve flagging. First, the prediction is
\emph{measurable}: the shrinkage law $\beta_{k}\sim 2/k$ is a specific
curve that truncated \SHA{} must trace if it is random-like, and a
reduced-round or structured map need not (\cref{sec:projects}). Second,
the practical fix is illuminating precisely because it \emph{breaks} the
composition. Standardized password-based key derivation does not compute
$\varphi^{k}$; it re-injects the password and a salt at every step,
iterating a map like $x \mapsto \SHA(x \shaxor \text{password})$ that
changes each round~\citep{NIST_SP800132}. That is no longer a single
function composed with itself---it is a \emph{driven} iteration, formally
the random walk of \cref{sec:prngs} rather than a functional graph---and
the salt-dependent driving is exactly what defeats the funnel (and the
precomputation attack of \cref{sec:itervuln}). The theory tells the
engineer why the naive design leaks and the standard one does not.

\subsection{The permutation at the bottom of the funnel}
\label{sec:permutation}

Where does the funnel end? Not at nothing: the image cannot shrink below
the points that lie on cycles, since a cyclic point is its own distant
iterate and survives forever. These \emph{recurrent} points---the set
$R=\bigcap_{k\ge 0}\varphi^{k}(\State)$, the union of all the cycles of
\cref{fig:rho}---are the funnel's floor, and here the user's intuition
lands exactly. \emph{On $R$, iterated \SHA{} is a genuine permutation.}
This is not a probabilistic near-miss; it is a theorem about every
finite self-map: restricted to its recurrent set a function is a
bijection, because each cyclic point has exactly one cyclic preimage
(the point before it on its cycle). \SHA{} permutes its cycles, and
$\varphi$ carried to the limit is a permutation of the roughly
$\sqrt{\pi N/2}\approx 2^{128.3}$ recurrent points---conjecturally a
\emph{uniformly random} permutation of a random such subset, which is
the strongest form of the random-map hypothesis.

So the imagined ``random permutation induced by iteration'' is real---%
but confined to an unreachable floor. The recurrent set is a
random-looking subset of size $2^{128.3}$, with no short description and
no way to enumerate it short of traversing orbits of length $2^{128}$;
one can prove it hosts a permutation without ever exhibiting a single
one of its elements. This is the same epistemic shape as the rest of the
paper (\cref{sec:telescope}): a structure we can characterize in full
and visit never.

Now the sharp contrast that answers ``why astronomically unlikely?''
Being a permutation on the \emph{recurrent set} is automatic; being a
permutation on the \emph{whole} space $\State$ is the rarest thing a map
can be. The fraction of self-maps that are permutations is
\[
  \frac{N!}{N^{N}} \;\approx\; \sqrt{2\pi N}\; e^{-N},
\]
so a uniformly random map is a permutation with probability of order
$e^{-N}$: not astronomically but \emph{doubly}-exponentially small, about
$2^{-1.44\cdot 2^{256}}$. Equivalently, injectivity anywhere is
expensive---a random map already has about $0.37 N$ points sharing their
images with others. Hence if \SHA{} were ever found to be injective on
its full domain, or even to cover meaningfully more than the predicted
$0.63N$ on the first step, that would not be a lucky coincidence: it
would be a gross structural signal, a deviation of the exact kind
\cref{sec:gap} is built to detect. The permutation lives at the bottom
of the funnel by necessity; a permutation at the \emph{top} would be a
five-alarm anomaly.

\subsection{``Unless\dots'': the one addition that breaks the bijection}
\label{sec:unless}

Why is \SHA{} not that anomaly---why does it funnel at all, rather than
permute? The answer is a single line of its construction, and it is worth
seeing because it identifies one-wayness with lossiness as the \emph{same
design act}. Recall from \cref{sec:roundshape} that the compression
function is a block cipher run in Davies--Meyer mode:
\[
  f(h, m) \;=\; E_{m}(h) \,\shaplus\, h,
\]
where $E_{m}$---the core of \SHA{}, essentially the cipher SHACAL-2---is,
for each fixed message block $m$, a \emph{permutation} of the state
(every round is invertible, so the whole cipher is a bijection one could
run backward). If $f$ were just $E_{m}$, iteration would be the iteration
of a permutation: reversible, lossless, its functional graph all cycles
and no trees, injective on the whole space---precisely the
doubly-exponentially unlikely object of \cref{sec:permutation}. The
feedforward ``$\shaplus\,h$'' is what destroys this
(\cref{fig:feedforward}). Adding the input back turns a bijection into a
generically many-to-one map---the sum of a permutation and the identity
need not be injective---and it is this one modular addition that makes
$f$ one-way, non-invertible, and lossy all at once. The tree-and-funnel
geometry of \cref{sec:fungraph,sec:shrinking}, the $37\%$ of leaves, the
$2/k$ shrinkage, the very existence of preimage-resistance: all of it is
manufactured by the single ``$\shaplus\,h$'' that the designers added
exactly so the rounds could not be run backward.

\begin{figure}[ht]
  \centering
  \begin{tikzpicture}[font=\small, >={Stealth[length=1.7mm]},
      pt/.style={circle, fill, inner sep=1.3pt}]
    % LEFT: E_m a permutation (bijection) — parallel arrows, no merging
    \begin{scope}
    \node[font=\scriptsize] at (0.9,1.35) {$E_{m}$ alone: bijection};
    \foreach \i in {0,1,2,3} {
      \node[pt] (l\i) at (0,-\i*0.5) {};
      \node[pt] (r\i) at (1.8,-\i*0.5) {};
    }
    \draw[->] (l0) -- (r1); \draw[->] (l1) -- (r3);
    \draw[->] (l2) -- (r0); \draw[->] (l3) -- (r2);
    \node[font=\scriptsize, align=center] at (0.9,-2.15)
      {every target hit once\\ \emph{reversible, lossless}};
    \end{scope}
    % RIGHT: f = E_m(h) + h — collisions and leaves appear
    \begin{scope}[xshift=5.4cm]
    \node[font=\scriptsize] at (0.9,1.35)
      {$f=E_{m}(h)\shaplus h$: collapses};
    \foreach \i in {0,1,2,3} {
      \node[pt] (a\i) at (0,-\i*0.5) {};
      \node[pt] (b\i) at (1.8,-\i*0.5) {};
    }
    \draw[->] (a0) -- (b1); \draw[->] (a1) -- (b1); % collision into b1
    \draw[->] (a2) -- (b3); \draw[->] (a3) -- (b3); % collision into b3
    \node[orange!80!black, font=\scriptsize, right] at (1.95,-0.5) {2 preimages};
    \node[orange!80!black, font=\scriptsize, right] at (1.95,-1.5) {2 preimages};
    \node[citewine, font=\scriptsize, right] at (1.95,0) {leaf};
    \node[citewine, font=\scriptsize, right] at (1.95,-1.0) {leaf};
    \node[font=\scriptsize, align=center] at (0.9,-2.15)
      {targets missed \& doubled\\ \emph{one-way, lossy}};
    \end{scope}
    \draw[->, thick] (2.55,-0.75) -- (3.35,-0.75)
      node[midway, above, font=\scriptsize] {$+\,h$};
  \end{tikzpicture}
  \caption{The ``unless.'' The cipher $E_{m}$ at the heart of \SHA{} is a
  permutation (left): iterate it and nothing is lost. The Davies--Meyer
  feedforward $f=E_{m}(h)\shaplus h$ (right) adds the input back, turning
  the bijection into a many-to-one map with leaves and collisions---the
  source of every tree, every lost bit, and all one-wayness. Iterated
  \SHA{} would be a random permutation on its whole domain
  (\cref{sec:permutation}) were it not for this one modular addition.}
  \label{fig:feedforward}
\end{figure}

This is the resolution of the opening puzzle. Iterated \SHA{} does induce
a random permutation---but on its recurrent floor, by necessity, not on
the whole space, which the feedforward forbids. The ``unless'' is
literal: \emph{unless one deletes the ``$\shaplus\,h$,''} in which case
\SHA{} degenerates into its underlying block cipher, iteration becomes
reversible, and the funnel disappears along with the security. The
lossiness the user intuited as a near-impossibility is, seen correctly,
the deliberate and load-bearing content of the design.

\subsection{Vulnerabilities specific to iteration}\label{sec:itervuln}

Iteration is not merely a lens; it is an attack surface and a building
block, and three concrete phenomena round out the picture.

\paragraph{Hash chains, and why they are only as strong as one step.}
Applying $\varphi$ repeatedly and revealing the chain \emph{backward} is
a classic authentication primitive: Lamport's one-time
passwords~\citep{Lamport1981} store $\varphi^{n}(s)$ on the server and
release $\varphi^{n-1}(s), \varphi^{n-2}(s), \dots$ one login at a time,
each value verified by one forward application and unforgeable ahead of
time by one-wayness. The same skeleton---iterate a one-way step, reveal
in reverse---underlies hash-based signatures (\cref{sec:uses}). What the
functional-graph view contributes is a caution: a chain is exactly as
strong as the one-wayness of a \emph{single} step, and if a short cycle
or a cheap partial inversion existed for the toy scales of
\cref{sec:projects}, the whole chain would unravel at once. Iteration
concentrates risk onto the per-step preimage problem.

\paragraph{Time--memory trade-offs: iteration as the attacker's
weapon.} The most striking vulnerability is that iterated hashing is the
data structure an adversary uses to \emph{invert} a one-way function by
precomputation. Hellman's time--memory trade-off~\citep{Hellman1980}
covers the domain with long chains $x, \varphi(x), \varphi^{2}(x),
\dots$, stores only the endpoints, and thereby inverts $\varphi$ on a
fresh target in online time $T$ and memory $M$ obeying
\[
  T\,M^{2} \;\approx\; N^{2},
\]
so that a one-time precomputation buys balanced online cost
$T = M = N^{2/3}$---far below the $N$ of blind brute force. Oechslin's
\emph{rainbow tables}~\citep{Oechslin2003} sharpen the constants with
per-column reduction functions and made the technique an everyday
password-cracking tool. At full \SHA{} scale $N^{2/3}=2^{170.7}$ is
safely unreachable, but against low-entropy inputs (short passwords) the
trade-off is devastating---and its defeat is once again the
\emph{driving} of \cref{sec:shrinking}: a per-user salt makes every
target's chain independent, so no shared precomputation
applies~\citep{NIST_SP800132}. The same iteration that stretches a key
for the defender builds the table for the attacker; salt is what breaks
the symmetry.

\paragraph{Fixed points and short cycles: cheap, sharp predictions.}
Finally, iteration poses questions with crisp random-model answers that
are eminently testable at toy scale. How many fixed points---strings
with $\SHA(x)=x$---should exist? For a random map the count is
$\mathrm{Poisson}(1)$: expected value exactly $1$, with probability
$e^{-1}\approx 0.37$ of \emph{none}. So \SHA{} very probably has a
handful (likely one) genuine $256$-bit fixed point that no one will ever
exhibit, a pleasant cousin of the ``does $\SHA(x)=x$ have a solution?''
puzzle. More generally the expected number of cycles of length $\ell$ is
$1/\ell$, a harmonic profile summing to the $\sqrt{\pi N/2}$ recurrent
points of \cref{sec:permutation}. Each of these is a falsifiable
prediction about a structureless map; a truncated \SHA{} that showed too
many fixed points, or cycles of anomalous length, would be exhibiting
structure. That nobody has looked, at model-organism scale, is the
recurring opportunity of this prospectus.

\begin{project}{The dynamics of iteration}{programming; a first
  probability course}{6--8 weeks (flagship)}
  Extend the functional-graph atlas of \cref{sec:telescope} from static
  statistics to \emph{dynamics}. For $\varphi_{n}(x) = $ first $n$ bits
  of $\SHA(x)$, $n = 16, \dots, 30$: (i) measure the image-shrinkage
  curve $\beta_{k}=|\varphi_{n}^{k}(\State)|/N$ and fit it against the
  $\beta_{k+1}=1-e^{-\beta_{k}}$ law and its $2/k$
  tail~\citep{FlajoletOdlyzko1990}; (ii) count fixed points and short
  cycles and compare with the $\mathrm{Poisson}(1)$ and $1/\ell$
  predictions; (iii) locate the recurrent set and confirm $\varphi$
  permutes it. Then repeat with reduced-round \SHA{} and with the
  feedforward \emph{removed} (pure $E_{m}$, a permutation) as a control,
  watching the funnel switch off. As a capstone, implement a small
  Hellman table~\citep{Hellman1980} on $\varphi_{n}$ and measure the
  $TM^{2}=N^{2}$ trade-off directly, with and without salt. Every
  quantity is exactly computable for $n \le 30$, and the ``remove the
  feedforward'' control makes the mechanism of \cref{sec:unless} visible
  as data.
\end{project}
